\subsection{An exact infinite-temperature endpoint criterion}
\label{sec:tripletendpoint}
The full triplet premise also removes the possible hidden interior cooling
interval: any anomaly begins at the high-temperature end of the Gibbs curve.
Only an endpoint derivative is needed to decide its existence.

\begin{corollary}[Cooling is decided by one endpoint covariance]
\label{cor:tripletendpoint}
Under the hypotheses and orientation of
Lemma~\ref{lem:firstthreearithmetic}, an ordinary asymptotic cooling effect
exists if and only if
\[
b_2'(0)<0,
\quad\text{equivalently}\quad
\operatorname{Cov}_{\mathrm{unif}}(w,E)>0.
\]
Here the uniform law is on all four states; $\beta=0$ is used only as an
analytic endpoint, not as an allowed witness preparation. If this condition
holds, there is exactly one finite positive stationary point of $b_2$.
It lies in $(0,\beta_b)$, and $|b_2|$ has its unique cooling maximum there.
The anomalous local ordering region is exactly $(0,\beta_{\rm turn})$.
If $b_2'(0)\geq0$, then $b_2'(\beta)>0$ for every finite $\beta>0$,
including when the endpoint derivative is zero, and no ordinary cooling
effect occurs. With a simple slow mode these conclusions apply to every
fixed norm asymptotically.
\end{corollary}
\begin{proof}
Write $b=b_2$. The contrast theorem gives $b<0$ on cooling and
$b'(\beta_b)>0$, as well as strict heating monotonicity. The derivative
at the bath is strictly positive directly: every prefix-contrast numerator
vanishes there with a strictly positive derivative, and the nonnegative
Perron coordinate vector is nonzero.

Put $C=w_3-w_4$. First suppose $C\geq0$. By the arithmetic restriction,
\[
m(0)=\frac{w_1+w_2+w_3}{3}\geq w_3\geq w_4,
\]
where $m$ is the conditional first-three Gibbs expectation from the sharp
bound. Since $m$ is nondecreasing, $b=Pm+(1-P)w_4$ gives
\[
b'=Pm'+P'(m-w_4)>0 \qquad (\beta\geq0).
\]
To justify strictness, if the first three components are not all equal,
the elementary contrast decomposition in the preceding proof gives
$m'>0$. Otherwise $m$ is constant, and nonconstancy of the full mode,
its zero bath mean and $w_4\leq0$ force $m>w_4$; then $P'>0$ supplies
strictness. Thus this case has no cooling anomaly.

Now suppose $C<0$. Use the derivative-curve representation
\eqref{eq:finitephi}: $b'=g\Phi$, where $g=\pi_1'>0$,
\[
\Phi(u)=A+Bu+C v(u),\qquad
u(\beta)\searrow-1 \text{ as }\beta\to\infty.
\]
The derivative-curve lemma makes $\Phi$ strictly concave in $u$.
It extends continuously to the two endpoints of
$[-1,u(0)]$, with
\[
\Phi(-1)=w_1-w_2\geq0,
\qquad \operatorname{sign}\Phi(u(0))=\operatorname{sign}b'(0).
\]
If $\Phi(u(0))\geq0$, strict concavity puts its value strictly above the
nonnegative chord at every interior point. Hence $b'(\beta)>0$ for every
finite $\beta>0$, even in the endpoint equality case.

If $\Phi(u(0))<0$, its positive value at the bath and its nonnegative
left endpoint imply that it has exactly one zero between $u(\beta_b)$
and $u(0)$. Indeed a concave function has an interval as its nonnegative
superlevel set; strict concavity and positivity at the bath rule out an
interior zero on the left-hand interval. On the right-hand interval the
slope of chords is decreasing, so after reaching zero the function cannot
return to nonnegative values. The zero is therefore unique, with negative
values to its right and positive values to its left. Since $u$ decreases
with $\beta$, $b'$ is negative for $0\leq\beta<\beta_{\rm turn}$ and
positive thereafter. The turning point lies strictly before the bath.
On cooling, $|b|=-b$, proving the stated maximum and anomalous interval.
Strict negativity persists at small finite positive inverse temperatures,
so the criterion supplies admissible finite witnesses.

Finally $b'(0)=-\operatorname{Cov}_{\mathrm{unif}}(w,E)$ by the fixed-mode
covariance identity. No interpretation of a stationary point alone is used.
\end{proof}

The covariance is explicitly
\[
\operatorname{Cov}_{\mathrm{unif}}(w,E)
=\frac14\sum_{i=1}^4w_iE_i
-\frac1{16}\Bigl(\sum_iw_i\Bigr)\Bigl(\sum_iE_i\Bigr).
\]
Thus within the triplet-failure family there is no need to scan the initial
cooling temperatures or solve an interior derivative equation merely to
decide existence. Section~\ref{sec:rateendpoint} gives an equivalent
evaluation using one cubic root and original-rate matrix products,
avoiding construction of the slow eigenvector. This is an asymptotic
criterion; it does not classify transient distance crossings.
